\section{Greedy Algorithm to Find Minimum Number of Coins}
\label{sec:minimum-coins}

\problemstatement{We are given a target amount $V$ and an unlimited supply of
coins (and notes) of standard denominations -- in the version used on the
sheet, the Indian currency denominations
\[
  \{1,\,2,\,5,\,10,\,20,\,50,\,100,\,500,\,1000\}.
\]
We must choose a multiset of these denominations whose values sum exactly to
$V$, using as \textbf{few coins/notes as possible}, and report that minimum
count (or the coins themselves).}

\begin{patternbox}
\emph{``Minimum number of coins/notes to make a given amount,''} when the
denomination set is a small, fixed, real-world currency system, is the
textbook greedy ``always use the largest denomination that fits'' problem.
The crucial, easy-to-miss keyword here is which \emph{specific} denomination
set is being used. If the problem instead says \emph{``given an arbitrary
list of coin denominations''} (as in the famous LeetCode ``Coin Change''
problem), the greedy approach can silently give a wrong answer, and you must
switch to dynamic programming. Recognizing \emph{which} of these two
near-identical-looking problems you are facing is itself the key skill being
tested here.
\end{patternbox}

\subsection*{Understanding the problem carefully}

For a fixed real-world currency system such as the Indian denominations
above, there is a special structural property: each denomination is large
enough, relative to the ones below it, that repeatedly taking the largest
denomination not exceeding the remaining amount never forces us into a
worse total count later. Concretely, every denomination in this list is at
least double the previous one except for the jumps $1 \to 2$, $2 \to 5$,
$5 \to 10$, $10 \to 20$, $20 \to 50$ -- and a short case check (which we will
not belabor here, since it is a property of \emph{this particular} currency
system rather than something we derive from first principles) confirms that
this set is a so-called \emph{canonical coin system}: the greedy algorithm
always produces an optimal (minimum-count) representation for it.

\begin{ideabox}[Greedy Choice]
Sort the denominations in decreasing order. Repeatedly subtract the largest
denomination that does not exceed the remaining amount, counting how many
times each denomination is used, until the remaining amount reaches zero.
\end{ideabox}

\subsection*{Why greedy works here, and why it can fail elsewhere}

The intuitive argument for why greedy is safe here is this: using one coin
of the largest available denomination $d$ removes as much value as possible
in a single coin. If using a smaller denomination instead could ever lead to
a strictly smaller total coin count, it would have to be because the smaller
coins ``line up'' more efficiently with the rest of the amount -- but because
each denomination in this system divides evenly into, or is closely
compatible with, the next one up, no such inefficiency can arise. This is
\emph{not} a universal truth about coin systems; it is a property that must
be verified for the specific set of denominations given.

To see the failure mode, suppose instead the denominations were
$\{1, 3, 4\}$ and the target amount were $V = 6$. Greedy would take one coin
of $4$, leaving $2$, then two coins of $1$, for a total of three coins:
$4+1+1=6$. But the optimal answer uses two coins of $3$: $3+3=6$, a total of
only two coins. Here greedy fails because $4$ is not a ``safe'' large
denomination relative to $3$ -- taking it first blocks the more efficient
combination. This is exactly why the general ``Coin Change: minimum coins''
problem (with an arbitrary denomination set) requires dynamic programming,
while this specific problem, restricted to real-world currency
denominations, is safe for greedy. Always check which version of the
problem you are being asked before reaching for greedy.

\subsection*{Algorithm}

\begin{enumerate}
  \item Store the denominations in decreasing order:
        $1000, 500, 100, 50, 20, 10, 5, 2, 1$.
  \item Initialize $\textit{remaining} \gets V$ and an empty result list.
  \item For each denomination $d$ in decreasing order:
  \begin{enumerate}
    \item While $d \le \textit{remaining}$: append $d$ to the result list,
          and set $\textit{remaining} \gets \textit{remaining} - d$.
  \end{enumerate}
  \item Return the result list (its length is the minimum coin count).
\end{enumerate}

\begin{algorithm}[H]
\caption{Minimum Number of Coins (Canonical System)}
\begin{algorithmic}[1]
\Procedure{MinimumCoins}{$V$}
  \State $\textit{denoms} \gets [1000, 500, 100, 50, 20, 10, 5, 2, 1]$
  \State $\textit{remaining} \gets V$
  \State $\textit{result} \gets$ empty list
  \For{each $d$ in \textit{denoms}}
    \While{$d \le \textit{remaining}$}
      \State append $d$ to \textit{result}
      \State $\textit{remaining} \gets \textit{remaining} - d$
    \EndWhile
  \EndFor
  \State \Return \textit{result}
\EndProcedure
\end{algorithmic}
\end{algorithm}

\subsection*{Dry run}

\dryrun{Take $V = 93$.}

\begin{itemize}
  \item $d = 1000$: $1000 > 93$, skip.
  \item $d = 500$: $500 > 93$, skip.
  \item $d = 100$: $100 > 93$, skip.
  \item $d = 50$: $50 \le 93$. Take one $50$. $\textit{remaining} = 43$.
        $50 \le 43$? No, stop this denomination.
  \item $d = 20$: $20 \le 43$. Take one $20$. $\textit{remaining} = 23$.
        $20 \le 23$. Take another $20$. $\textit{remaining} = 3$.
        $20 \le 3$? No, stop.
  \item $d = 10$: $10 > 3$, skip.
  \item $d = 5$: $5 > 3$, skip.
  \item $d = 2$: $2 \le 3$. Take one $2$. $\textit{remaining} = 1$.
        $2 \le 1$? No, stop.
  \item $d = 1$: $1 \le 1$. Take one $1$. $\textit{remaining} = 0$.
\end{itemize}

The coins used are $\{50, 20, 20, 2, 1\}$, a total of $5$ coins, and indeed
$50 + 20 + 20 + 2 + 1 = 93$. No combination of these denominations can make
$93$ with fewer than $5$ coins.

\subsection*{C++ implementation}

\begin{lstlisting}
#include <bits/stdc++.h>
using namespace std;

vector<int> minimumCoins(int V) {
    vector<int> denoms = {1000, 500, 100, 50, 20, 10, 5, 2, 1};
    vector<int> result;

    int remaining = V;
    for (int d : denoms) {
        while (d <= remaining) {
            result.push_back(d);
            remaining -= d;
        }
    }
    return result;
}

int main() {
    int V = 93;
    vector<int> coins = minimumCoins(V);

    cout << "Number of coins used: " << coins.size() << endl;
    cout << "Coins: ";
    for (int c : coins) cout << c << " ";
    cout << endl;
    return 0;
}
\end{lstlisting}

\begin{complexitybox}
\textbf{Time Complexity.} There are only $9$ fixed denominations, so the
outer loop is $O(1)$ in the number of denominations. The inner
\textbf{while} loop, across all denominations, executes at most once per
coin actually placed into the result, and the number of coins used is
$O(\log V)$ in the worst case for this denomination system (since
denominations grow roughly geometrically). So the overall time complexity is
\[
  O(\log V),
\]
effectively a constant for realistic input sizes.
\textbf{Space Complexity.} We store the result list of coins used, which has
length $O(\log V)$, so the space complexity is $O(\log V)$ for the output
(or $O(1)$ auxiliary space if we only need the \emph{count} rather than the
list of coins itself).
\end{complexitybox}

\begin{takeawaybox}
Greedy coin selection is only correct for \emph{canonical} denomination
systems -- ones where the structure of the denominations guarantees that
taking the largest coin first never blocks an optimal solution. Always
distinguish this fixed-currency version of the problem from the general
``minimum coins from an arbitrary denomination list'' problem, which has no
such guarantee and requires dynamic programming (typically an unbounded
knapsack-style DP over amounts $0, \dots, V$).
\end{takeawaybox}
